\scanpage{254}
\mspage{1\deleted{\unsure{3}}2\deleted{\unsure{9|8}}8}
% archive sheet number: 253

Considérons maintenant le diagramme \inserted{(traits pleins)}

\diagram{\begin{tikzcd}[ampersand replacement=\&,row sep=large,column sep=small]
\deleted{X\amalg X}\arrow[d]\arrow[d,decorate,decoration={zigzag,segment length=3pt,amplitude=2pt},no head]
 \& X\vee X\arrow[dl,"{(\partial_0^X,\partial_1^X)}"']\arrow[dr,"{(\partial_0'{}^X,\partial_{\unsure{1|0}}'{}^X)}"] \& {} \\
I\times X\arrow[rr,"{\unsure{\varphi|\phi}\times\mathrm{id}_X}","{\substack{\|\\\deleted{\unsure{u}}\varphi_X}}"' ]
 \arrow[dr,"h"']\arrow[d,"\deleted{h}"']\arrow[d,decorate,decoration={zigzag,segment length=3pt,amplitude=3pt},no head]
 \& {} \& \deleted{\unsure{I}}J\times X\arrow[dl,dashed,"h'"] \\
{} \& Y \& {}
\end{tikzcd}}

, où $\partial'_0{}^X,\partial'_1{}^X$ sont définis comme $\partial_0^X,\partial_1^{\unsure{\alpha|X}}$, mais utilisant $j_0,j_1\in\mathrm{Ob}\,J$ au lieu de $i_0,i_1\in\mathrm{Ob}\,I$. S'il existait dans $\mathrm{Cat}$ une flèche $h'$ (pointillée) telle que $h'\circ\deleted{\unsure{u}}\varphi_X=h$, il s'ensuivrait que \deleted{les flèches} $h'$ est une homotopie de $f_0$ et $f_1$ suivant $(J,j_0,j_1)$.

Comme $J$ est \unsure{$0$-connexe|$o$-connexe}, il s'ensuivrait aussitôt que $f_0,f_1$ sont homotopes au sens habituel, dans $\mathrm{Cat}$, donc \deleted{\unsure{f}} que $f_0$ et $f_1$ ont $=$ image dans $\mathrm{Hot}_W$. Mais si on ne peut trouver une flèche $h'$ dans $\mathrm{Cat}$, du moins peut-on la trouver dans $\mathrm{Hot}_W$, puisque $\varphi_X\in W$ donc un iso dans $\mathrm{Hot}_W$ (et $h'$ dans $\mathrm{Hom}_W(J\times X,Y)$ \unsure{sera|en.} $=$ \unsure{uniquement|univoquement} déterminée). Donc dans $\mathrm{Hot}_W$ on a \deleted{\unsure{comme}}
\[
 f_0=h'\partial'_0{}^X,\qquad f_1=h'\partial'_1{}^X.
\]
